\subsection{A common scalar history with readout and restoration}
\label{sec:common-history-feedback}

The finite golden scalar action admits an operational composition in which
readout changes the field and retained records restore it before subsequent
evolution. This connects the local instrument to the same history whose
action and separated response are checked. It strengthens a free-field
execution by accounting for the later consequences of readout error.
The construction uses the declared $64$-site operator $A$, masses $M$,
left-supported velocity $v$, right-supported detector $g$, and
$\tau^2=1/245$. It retains $\|v\|_M\le3/4$, $\|g\|_M\le1/4$,
$\sum_iM_i\le1$, and $0<\tau^2 A<3I$.
The source and detector supports have coordinate separation greater than
$3/8$. These are finite supplied coordinates and operations.

At each detector site a zero pointer and the updated field value
$x$ undergo the classical averaging probe
$(x,0)\mapsto(x/2,x/2)$. A retained response
$r=x/2+\eta$ is routed to the readout and used to restore the field to
$2r$ and reset the pointer. Consequently the actual subsequent action
consumes $x+2\eta$, rather than an independently reloaded reference field.
Every probe, retained response, feedback, route, accumulator reset and
clock update names its actual consumed writers and values.

\begin{proposition}[Accumulated readout error on one action history]
\label{prop:common-history-feedback-error}
Let $\tau>0$, let $A$ be positive and self-adjoint in the $M$ inner product,
with $0\le\tau^2 A\le4I$, and let $\sum_iM_i\le1$.
For nonnegative error budgets $\epsilon_p,\epsilon_r$, suppose
the initial velocity differs by at most $\epsilon_p$ in $M$ norm,
and every retained local response error obeys
$|\eta_{j,i}|\le\epsilon_r$. The restored history satisfies
\[
 \widetilde u_j=(2I-\tau^2A)\widetilde u_{j-1}
                 -\widetilde u_{j-2}+2\eta_j,
 \qquad
 \|\widetilde u_j-u_j\|_M
 \le j\tau\epsilon_p+j(j+1)\epsilon_r,
\]
where $u_j$ is the ideal action history, both histories start at zero
position, and $\eta_j$ vanishes outside the detector support.
\end{proposition}
\begin{proof}
Restoration has exact error $2\eta_j$, with norm at most
$2\epsilon_r$. In a spectral mode put
$c=1-\tau^2\lambda/2\in[-1,1]$. The response to an additive impulse
after $n$ further steps is the Chebyshev polynomial $U_n(c)$.
Its representation $U_n(\cos\theta)=\sum_{k=0}^n
e^{i(n-2k)\theta}$ gives $|U_n(c)|\le n+1$, including the
endpoints by continuity. Spectral calculus and the triangle inequality
therefore give $j\tau\epsilon_p+
2\epsilon_r\sum_{k=1}^j k$.
\end{proof}

For two independently perturbed baseline and intervention histories,
the corresponding detector error at step $j$ is bounded by
\[
 2\|g\|_M\{j\tau\epsilon_p+j(j+1)\epsilon_r\}
       +2\epsilon_d,
\]
where $\epsilon_d$ bounds each final scalar readout error.
No independence or probabilistic cancellation of the errors is assumed.
Routing, reset, averaging, restoration arithmetic and accumulation are
exact in this contract. Errors at those additional operation locations
would require their own propagation budget.
Take $\epsilon_p=\epsilon_r=10^{-6}$,
$\epsilon_d=10^{-7}$ and model-clock uncertainty $10^{-6}$.
The nominal action-step clock selects the last completed record over
$[15\tau,17\tau]$; it consumes no future observation.
The clock uncertainty bounds displacement of the comparison time while
holding this nominal record selection fixed.
For the continuous-time reference on the same finite action,
$U(t)=A^{-1/2}\sin(t\sqrt A)v$,
\[
 |D'(U(t))|\le\|g\|_M\|v\|_M\le3/16.
\]
Thus holding and clock uncertainty contribute at most
$(3/16)(\tau+10^{-6})$. The full finite-action discretization error
is the independently reconstructed all-mode bound at the retained
steps, multiplied by $\|g\|_M$. Preparation, propagated feedback,
final readout, discretization, hold reconstruction and timing errors
are added before comparison with the response. Exact algebraic
arithmetic and exact finite supports contribute zero numerical and
localization error on this declared model; this is not a spatial
continuum error assertion.

The independently replayed packet resolves a continuous model-time
window of length $2\tau>0.127$. Its conservative total comparison
error is below $0.01364$, while every admitted paired held readout
has magnitude above $0.01576$. The strict positive margin also bounds
the continuous same-action signal away from zero across that window.
The explicit noisy histories are examples inside the uniform bound.
A doubled-input history checks linear response. Disabling restoration
produces a different field history and breaks the identification of the
decoded records with the evolved state. The consumer rejects resealed
substitutions of the action, population, clock, readout, record value
or consumed writer.

This is a classical software realization with bounded registers at the
declared cutoff, protected coefficients, supplied record channels and
exact retained algebraic values. All auxiliary events and coefficient
precision are counted separately from the model-action clock; their
physical durations and costs are not inferred from that clock. Quantum
instrument semantics, physical source selection and laboratory
calibration are separate obligations. The replay is in
\path{code/common_history_packet/verify.py}; its complete histories and
comparison certificate are in
\path{evidence/common_history_packet_20260925/}.
The four declarations in
\path{Lean/Screen/CommonHistoryFeedback.lean}
verify exact restoration and its finite accumulation factor. The
operator stability and spectral error argument above are analytical.
